% 07_defenses.tex — defense blueprint + matrix.

\section{Defense Strategies}\label{sec:defenses}

The RQ-4 finding --- that the composition of four defenses
reduces leak rate to zero, but no single defense closes the gap
--- motivates a layered defense blueprint. We enumerate the
four defenses, their cost, and the boundary they address.

\subsection{Defense matrix}

\begin{table}[t]
\caption{Defense matrix. Each row is one defense; columns are
the boundary it addresses, the mechanism, and the deployment
cost (Low = configuration change; Medium = middleware
addition; High = protocol-level rewrite).}
\label{tab:defense-matrix}
\centering
\small
\begin{tabularx}{\linewidth}{l l X c}
\toprule
Defense & Boundary & Mechanism & Cost \\
\midrule
Per-tenant tool registry         & namespace    & Static allow-list of tools per tenant                              & Low \\
Tenant-prefixed cache keys      & cache        & Cache keys include \texttt{tenant\_id}                             & Low \\
Resource-path canonicalisation  & resource     & Reject URIs that escape tenant root after \texttt{realpath}        & Low \\
Audience-bound JWT tokens       & auth         & \texttt{aud} claim enforced at every await boundary               & Medium \\
$m$TLS / channel auth           & transport    & Per-pipe tenant identity via SAN pinning                           & Medium \\
\bottomrule
\end{tabularx}
\end{table}

\subsection{Per-defense specification}

\paragraph{Per-tenant tool registry (namespace).} The
\texttt{Defenses.per\_tenant\_tool\_registry} flag toggles a
middleware that filters \texttt{tools/list} and
\texttt{tools/call} responses by the requesting tenant. The
registry is a \texttt{dict[tenant\_id, set[tool\_name]]}
populated at server startup; an unknown tenant receives an
empty list, an unknown tool receives a 403. The middleware
adds a constant-time hash lookup per call ($\sim 50\,\mu s$ in
our reference implementation).

\paragraph{Tenant-prefixed cache keys (cache).} The
\texttt{Defenses.tenant\_prefixed\_cache\_keys} flag
re-computes every cache key as
\texttt{(tenant\_id, tool\_name, args\_hash)} instead of
\texttt{(tool\_name, args\_hash)}. The middleware is a single
line of code at the cache lookup site; cost is negligible.

\paragraph{Resource-path canonicalisation (resource).} The
\texttt{Defenses.resource\_path\_canonicalisation} flag
replaces the resolver's string-prefix check with
\texttt{os.path.realpath} after decoding percent-encoded
segments. URIs that escape the tenant root return a 404. The
middleware adds $\sim 200\,\mu s$ per
\texttt{resources/read} call.

\paragraph{Audience-bound JWT tokens (auth).} The
\texttt{Defenses.mtls} flag (the prompt-mandated name for the
audience-bound JWT defense) issues JSON Web Tokens with an
\texttt{aud} claim scoped to the tenant identifier. The
server re-validates the \texttt{aud} claim at every await
boundary (re-check, not single-check, to defeat the
TOCTOU race in misuse case). The middleware adds $\sim
300\,\mu s$ per call.

\subsection{Defense-composition rationale}

No single defense is sufficient because each addresses a
single boundary:

\begin{itemize}
  \item The per-tenant tool registry does not affect cache
        poisoning (RQ-2 demonstrates cache attacks bypass
        the registry).
  \item Tenant-prefixed cache keys do not affect resource
        traversal (RQ-3 demonstrates resource attacks bypass
        the cache middleware).
  \item Resource canonicalisation does not affect prompt
        injection (the prompt is not a URI).
  \item Audience-bound JWTs do not affect cache key
        collisions (the collision is below the auth layer).
\end{itemize}

The composition closes the gap because each attack's
exploitation path crosses \emph{at least one} boundary whose
defense is engaged. Cache attacks cross the namespace boundary
(tool routing) and the cache boundary; the registry defence
catches the former and the cache defence catches the latter.
Prompt-injection attacks cross the resource boundary (the
injected content is read via \texttt{resources/read}) and the
auth boundary (the injected instruction seeks elevation); the
canonicalisation and JWT defences catch them in series.

\subsection{Layered defense blueprint}\label{sec:blueprint}

We recommend the following order of adoption for a production
MCP deployment:

\begin{enumerate}
  \item \textbf{Audience-bound JWTs (auth).} Highest leverage;
        addresses the broadest set of attacks
        (token-replay, capability-spoofing, SSRF).
  \item \textbf{Per-tenant tool registry (namespace).} Low
        cost; addresses tool-squatting and confused-deputy
        attacks.
  \item \textbf{Tenant-prefixed cache keys (cache).} Low cost;
        addresses the dominant leakage surface (RQ-2).
  \item \textbf{Resource-path canonicalisation (resource).} Low
        cost; addresses path traversal and symlink escape.
\end{enumerate}

The order is chosen so that each subsequent defense is
additive rather than overlapping: the JWT defence gates
\emph{who}, the tool registry gates \emph{what}, the cache
keys gate \emph{how often}, and the resource canonicalisation
gates \emph{where to}.

\subsection{Defense overhead (measured)}

The per-defense overhead is measured on the in-process
connector by \texttt{analysis/scripts/bench\_defenses.py}
(10{,}000 iterations per defense, median reported):

\begin{itemize}
  \item Per-tenant tool registry: $0.1\,\mu s$ per call
        (constant-time hash lookup).
  \item Tenant-prefixed cache keys: $0.1\,\mu s$ per call
        (single string concatenation).
  \item Resource-path canonicalisation: $\sim 113\,\mu s$ per
        call (\texttt{os.path.realpath} on a typical URI).
  \item Audience-bound JWTs: $\sim 2\,\mu s$ per call
        (HMAC-SHA256 verify).
\end{itemize}

The cumulative per-call overhead is therefore
$\sim 115\,\mu s$ on the in-process connector. On a real
HTTP+SSE transport the network round-trip ($\sim 1$--$10\,ms$
locally) dominates the per-defense overhead; the relative
ordering of the defenses (and the composition result) is
unaffected. The benchmark harness is included in the
artifact bundle so reviewers can re-measure on their target
transport.

\subsection{Defense failure modes}

Each defense has a failure mode that a determined attacker can
exploit. We document them so that adopters understand the
defenses' residual risk:

\paragraph{Per-tenant tool registry failure mode.} The
registry is a \texttt{dict[tenant\_id, set[tool\_name]]};
if the attacker can register a tool with a \texttt{tenant\_id}
the server trusts (e.g. via a privilege-escalation bug in the
host's authentication), the registry's allow-list is moot.
The defense therefore relies on the host's authentication
chain, which is itself defended by the audience-bound JWT
defense (the second layer).

\paragraph{Tenant-prefixed cache keys failure mode.} If the
attacker can poison a tenant's cache key (e.g. via a separate
write-side bug), the prefix does not prevent the poisoned
key from being read by the same tenant. The defense therefore
relies on the write-side defenses (per-tenant tool registry
and resource canonicalisation) to prevent the initial
poisoning. The defense closes the \emph{read-side} leakage
path, not the \emph{write-side} attack.

\paragraph{Resource-path canonicalisation failure mode.} If
the attacker can craft a URI that resolves within the tenant
root but contains attacker content from another tenant (e.g.
via a server-side include or a symlink planted at startup),
the canonicalisation does not detect the cross-tenant
content. The defense closes the \emph{path-traversal}
leakage path, not the \emph{content-injection} attack. The
prompt-injection defense (future work) would close the
latter.

\paragraph{Audience-bound JWTs failure mode.} If the attacker
can compromise the upstream identity provider (out of scope
in our threat model), the JWT's \texttt{aud} claim is moot
because the attacker can mint a token with any
\texttt{aud}. The defense relies on the identity provider's
security; the protocol-layer defense is the second line of
defence, not the first.

\subsection{Defenses and the four misuse cases}

We map each defense to the misuse cases it addresses:

\begin{itemize}
  \item \textbf{MC-1 (tool shadowing).} Closed by the
        per-tenant tool registry. The shadowed \texttt{set\_env}
        tool is not in Tenant~B's allow-list; the call
        returns a 403.
  \item \textbf{MC-2 (session fixation).} Closed by the
        audience-bound JWTs. The attacker cannot mint a
        session-resumption JWT with Tenant~B's
        \texttt{aud} claim.
  \item \textbf{MC-3 (resource path traversal).} Closed by
        the resource-path canonicalisation. The
        percent-encoded slashes are decoded and
        \texttt{realpath}-ed before the tenant-prefix check;
        the resolver returns a 404.
  \item \textbf{MC-4 (embedding cache poisoning).} Closed by
        the tenant-prefixed cache keys. The cache key now
        includes \texttt{tenant\_id}; a collision crafted by
        Tenant~A does not match Tenant~B's key.
\end{itemize}

The composition is therefore complete: every documented
misuse case is closed by at least one defense, and most are
closed by two (defense-in-depth). The RQ-4 empirical result
($L = 0$ at composition) is the numerical confirmation of
this mapping.